\begin{exercise}{79}
Show that an inner fibration $\mathcal C\to\mathcal D$ between $\infty$-categories is
an isofibration if and only if the induced functor
$\mathrm N(h\mathcal C)\to\mathrm N(h\mathcal D)$ is an isofibration.

\begin{proof}
For any ordinary functor $q\colon A\to B$, its nerve is an
inner fibration: fill the top inner horn in $\mathrm N(A)$;
uniqueness of the filler in $\mathrm N(B)$ forces its image
to be the prescribed bottom simplex (Theorem~1.1.52).
By Corollary~1.2.14 and Definition~1.2.22, equivalences in
ordinary nerves are exactly isomorphisms.
Definition~2.1.4 therefore says that $\mathrm N(q)$ is an
isofibration exactly when every isomorphism
$\beta\colon q(a)\to b$ lifts to an isomorphism
$\alpha\colon a\to a'$ with $q(a')=b$ and $q(\alpha)=\beta$.
Call this the ordinary isofibration condition on $q$.

Write $p\colon\mathcal C\to\mathcal D$ for the given inner
fibration. Suppose first that $p$ is an isofibration.
For $x\in\mathcal C$ and an isomorphism
$\beta\colon p(x)\to b$ in $h\mathcal D$,
Corollary~1.2.13 gives a representative edge
$f\colon p(x)\to b$ in $\mathcal D$.
It is an equivalence by Definition~1.2.22.
Definition~2.1.4 lifts $f$ to an equivalence
$e\colon x\to y$ with $p(e)=f$.
Then $[e]$ is an isomorphism lifting $\beta$ under $hp$.
Thus $hp$ satisfies the ordinary isofibration condition,
so $\mathrm N(hp)$ is an isofibration.

Conversely, suppose that $\mathrm N(hp)$ is an isofibration.
Given $x\in\mathcal C$ and an equivalence
$f\colon p(x)\to b$ in $\mathcal D$, the ordinary
isofibration condition gives an isomorphism
$\alpha\colon x\to y$ in $h\mathcal C$ with
$p(y)=b$ and $hp(\alpha)=[f]$.
Choose an edge $a\colon x\to y$ representing $\alpha$.
The equality $[p(a)]=[f]$, together with Corollary~1.2.13
and Definition~1.2.3, gives a $2$-simplex $\theta$ in
$\mathcal D$ with
\[
 d_2\theta=p(a),\qquad d_1\theta=f,\qquad d_0\theta=s_0(b).
\]
Its inner horn lifts along $p$ with prescribed edges
$a$ and $s_0(y)$, since their images are $d_2\theta$
and $d_0\theta$.
The long edge $e\colon x\to y$ of the lifted triangle
satisfies $p(e)=f$ and
$[e]=[s_0(y)]\circ[a]=\alpha$.
Hence $e$ is an equivalence and solves the required lifting
problem. This proves that $p$ is an isofibration.
\end{proof}
\end{exercise}
